\section{USAJMO 2014}
        \begin{enumerate}
        	\item (\textit{USAJMO 2014}) Let $a$, $b$, $c$ be real numbers greater than or equal to $1$. Prove that 
 \[\min{\left (\frac{10a^2-5a+1}{b^2-5b+10},\frac{10b^2-5b+1}{c^2-5c+10},\frac{10c^2-5c+1}{a^2-5a+10}\right )}\leq abc.\]


 



\textbf{Solution}

Since $(a-1)^5\ge 0$,
 \[a^5-5a^4+10a^3-10a^2+5a-1\ge 0\]or
 \[10a^2-5a+1\le a^3(a^2-5a+10)\]Since $a^2-5a+10=\left( a-\dfrac{5}{2}\right)^2 +\dfrac{15}{4}>0$,
 \[\frac{10a^2-5a+1}{a^2-5a+10}\le a^3\]Also note that $10a^2-5a+1=10\left( a-\dfrac{1}{4}\right)^2+\dfrac{3}{8}> 0$,We conclude
 \[0\le \frac{10a^2-5a+1}{a^2-5a+10}\le a^3\]Similarly, 
 \[0\le \frac{10b^2-5b+1}{b^2-5b+10}\le b^3\]
 \[0\le \frac{10c^2-5c+1}{c^2-5c+10}\le c^3\]So 
 \[\left(\frac{10a^2-5a+1}{a^2-5a+10}\right)\left(\frac{10b^2-5b+1}{b^2-5b+10}\right)\left(\frac{10c^2-5c+1}{c^2-5c+10}\right)\le a^3b^3c^3\]or
 \[\left(\frac{10a^2-5a+1}{b^2-5b+10}\right)\left(\frac{10b^2-5b+1}{c^2-5c+10}\right)\left(\frac{10c^2-5c+1}{a^2-5a+10}\right) \le(abc)^3\]Therefore,
 \[\min\left(\frac{10a^2-5a+1}{b^2-5b+10},\frac{10b^2-5b+1}{c^2-5c+10},\frac{10c^2-5c+1}{a^2-5a+10}\right )\le abc.\]

 


	\item (\textit{USAJMO 2014}) Let $\triangle{ABC}$ be a non-equilateral, acute triangle with $\angle A=60^\circ$, and let $O$ and $H$ denote the circumcenter and orthocenter of $\triangle{ABC}$, respectively.


 (a) Prove that line $OH$ intersects both segments $AB$ and $AC$.


 (b) Line $OH$ intersects segments $AB$ and $AC$ at $P$ and $Q$, respectively. Denote by $s$ and $t$ the respective areas of triangle $APQ$ and quadrilateral $BPQC$. Determine the range of possible values for $s/t$.


 



\textbf{Solution}


                    \begin{figure}[H]
                        \centering
                        \includegraphics[width=0.5\textwidth]{Solution/USAJMO_I_2014_2_Solution_Figure_1.png}
                    \end{figure}
                    

 Lemma: $H$ is the reflection of $O$ over the angle bisector of $\angle BAC$ (henceforth 'the' reflection)

 Proof: Let $H'$ be the reflection of $O$, and let $B'$ be the reflection of $B$. 

 Then reflection takes $\angle ABH'$ to $\angle AB'O$.

 $\Delta ABB'$ is equilateral, and $O$ lies on the perpendicular bisector of $\overline{AB}$ 

 It's well known that $O$ lies strictly inside $\Delta ABC$ (since it's acute), meaning that $\angle ABH' = \angle AB'O = 30^{\circ},$ from which it follows that $\overline{BH'} \perp \overline{AC}$ . Similarly,  $\overline{CH'} \perp \overline{AB}$. Since $H'$ lies on two altitudes, $H'$ is the orthocenter, as desired.

 End Lemma

 So $\overline{OH}$ is perpendicular to the angle bisector of $\angle OAH$, which is the same line as the angle bisector of $\angle BAC$, meaning that $\Delta APQ$ is equilateral. 

 Let its side length be $s$, and let $PH=t$, where $0 < t < s, t \neq s/2$ because $O$ lies strictly within $\angle BAC$, as must $H$, the reflection of $O$. Also, it's easy to show that if $O=H$ in a general triangle, it's equilateral, and we know $\Delta ABC$ is not equilateral. Hence H is not on the bisector of $\angle BAC \implies t \neq s/2$. Let $\overrightarrow{BH}$ intersect $\overline{AC}$ at $P_B$. 

 Since $\Delta HP_BQ$ and $BP_BA$ are 30-60-90 triangles, $AB=2AP_B=2(s-QP_B)=2(s-HQ/2)=2s-HQ=2s-(s-t)=s+t$

 Similarly, $AC=2s-t$

 The ratio $\frac{[APQ]}{[ABC]-[APQ]}$ is $\frac{AP \cdot AQ}{AB \cdot AC - AP \cdot AQ} = \frac{s^2}{(s+t)(2s-t)-s^2}$The denominator equals $(1.5s)^2-(.5s-t)^2-s^2$ where $.5s-t$ can equal any value in $(-.5s, .5s)$ except $0$. Therefore, the denominator can equal any value in $(s^2, 5s^2/4)$, and the ratio is any value in $\boxed{\left(\frac{4}{5},1\right)}.$

 Note: It's easy to show that for any point $H$ on $\overline{PQ}$ except the midpoint, Points B and C can be validly defined to make an acute, non-equilateral triangle.

 



\textbf{Solution 2}

Let $J$ be the farthest point on the circumcircle of $\triangle{ABC}$ from line $BC$.Lemma: Line $AJ$||Line $OH$Proof: Set $b=-\dfrac{1}{2}+i\dfrac{\sqrt{3}}{2}$ and $c=-\dfrac{1}{2}-i\dfrac{\sqrt{3}}{2}$, and $a$ on the unit circle. It is well known that $o=0$ and $h=a+b+c$, so we have $h-o=a+b+c=a-1=a-j$, so $\dfrac{a-j}{h-o}$ is real and thus the 2 lines are parallel.

 WLOG let $A$ be in the first quadrant. Clearly by the above lemma $OH$ must intersect line $BC$ closer to $B$ than to $C$. Intersect $AJ$ and $BC$ at $D$ and $OH$ and $BC$ at $E$. We clearly have $0 < \dfrac{\overarc{JC}-\overarc{AB}}{2} = \dfrac{120-\overarc{AC}}{2} = \angle{ADB} = \angle{OEC} < 30 = \angle{OBC}$, $OH$ must intersect $AB$. We also have, letting the intersection of line $OH$ and line $AC$ be $Q$, and letting intersection of $OH$ and $AB$ be $P$, $\angle{OQA} = \angle{JAB} = \dfrac{\overarc{JB}}{2} = 60$. Since $\angle{OCA}=\angle{BCA}-\angle{BCO} < 90-30 =\angle{OQA}$, and $\angle{OQC} = 120>\angle{OAC}$, $OH$ also intersects $AC$. We have $\angle{OPA}=\angle{PAJ}=60$, so $\triangle{APQ}$ is equilateral. Letting $AJ=2x$, and letting the foot of the perpendicular from $O$ to $AJ$ be $L$, we have $OL=\sqrt{1-x^{2}}$, and since $OL$ is an altitude of $\triangle{APQ}$, we have $[APQ]=\dfrac{OL^{2}\sqrt{3}}{3} = \dfrac{\sqrt{3}(1-x^{2})}{3}$. Letting the foot of the perpendicular from $A$ to $OJ$ be $K$, we have $\triangle{JKA}\sim \triangle{JLO}$ by AA with ratio $\dfrac{JA}{JO} = 2x$. Therefore, $JK = 2x(JL) = 2x^{2}$. Letting $D$ be the foot of the altitude from $J$ to $BC$, we have $KD=JD-JK=\dfrac{3}{2}-2x^{2}$, since $re(B)=re(C) \implies JD=re(j)-(-\dfrac{1}{2})=\dfrac{3}{2}$. Thus, since $BC=\sqrt{3}$ we have $[ABC]=\dfrac{(\dfrac{3}{2}-2x^{2})(\sqrt{3})}{2} = \dfrac{(3-4x^{2})(\sqrt{3})}{2}$, so $[PQCB]=[ABC]-[APQ]=\dfrac{(5-8x^{2})(\sqrt{3})}{12}$, so $\dfrac{[APQ]}{[PQCB]} = \dfrac{4-4x^{2}}{5-8x^{2}} = \dfrac{1}{2}+\dfrac{3}{10-16x^{2}}$. We have $x=\sin(\dfrac{JOA}{2})$, with $0<120-2\angle{ACB}=\angle{JOA}<60$, so $x$ can be anything in the interval $(0, \dfrac{1}{2})$. Therefore, the desired range is $(\dfrac{4}{5}, 1)$.

 Solution by Shaddoll

 


	\item (\textit{USAJMO 2014}) Let $\mathbb{Z}$ be the set of integers. Find all functions $f : \mathbb{Z} \rightarrow \mathbb{Z}$ such that 
 \[xf(2f(y)-x)+y^2f(2x-f(y))=\frac{f(x)^2}{x}+f(yf(y))\] for all $x, y \in \mathbb{Z}$ with $x \neq 0$.


 



\textbf{Solution}

Note: This solution is kind of rough. I didn't want to put my 7-page solution all over again. It would be nice if someone could edit in the details of the expansions.

 Lemma 1: $f(0) = 0$.Proof: Assume the opposite for a contradiction. Plug in $x = 2f(0)$ (because we assumed that $f(0) \neq 0$), $y = 0$. What you get eventually reduces to:
 \[4f(0)-2 = \left( \frac{f(2f(0))}{f(0)} \right)^2\]which is a contradiction since the LHS is divisible by 2 but not 4.

 Then plug in $y = 0$ into the original equation and simplify by Lemma 1. We get:
 \[x^2f(-x) = f(x)^2\]Then:

 
 \begin{align*} x^6f(x) &= x^4\bigl(x^2f(x)\bigr)\\ &= x^4\bigl((-x)^2f(-(-x))\bigr)\\ &= x^4(-x)^2f(-(-x))\\ &= x^4f(-x)^2\\ &= f(x)^4 \end{align*}


 Therefore, $f(x)$ must be 0 or $x^2$.

 Now either $f(x)$ is $x^2$ for all $x$ or there exists $a \neq 0$ such that $f(a)=0$. The first case gives a valid solution. In the second case, we let $y = a$ in the original equation and simplify to get:
 \[xf(-x) + a^2f(2x) = \frac{f(x)^2}{x}\]But we know that $xf(-x) = \frac{f(x)^2}{x}$, so:
 \[a^2f(2x) = 0\]Since $a$ is not 0, $f(2x)$ is 0 for all $x$ (including 0). Now either $f(x)$ is 0 for all $x$, or there exists some $m \neq 0$ such that $f(m) = m^2$. Then $m$ must be odd. We can let $x = 2k$ in the original equation, and since $f(2x)$ is 0 for all $x$, stuff cancels and we get:
 \[y^2f(4k - f(y)) = f(yf(y))\]\textbf{for $\mathbf{k \neq 0}$.}Now, let $y = m$ and we get:
 \[m^2f(4k - m^2) = f(m^3)\]Now, either both sides are 0 or both are equal to $m^6$. If both are $m^6$ then:
 \[m^2(4k - m^2)^2 = m^6\]which simplifies to:
 \[4k - m^2 = \pm m^2\]Since $k \neq 0$ and $m$ is odd, both cases are impossible, so we must have:
 \[m^2f(4k - m^2) = f(m^3) = 0\]Then we can let $k$ be anything except 0, and get $f(x)$ is 0 for all $x \equiv 3 \pmod{4}$ except $-m^2$.  Also since $x^2f(-x) = f(x)^2$, we have $f(x) = 0 \Rightarrow f(-x) = 0$, so $f(x)$ is 0 for all $x \equiv 1 \pmod{4}$ except $m^2$. So $f(x)$ is 0 for all $x$ except $\pm m^2$. Since $f(m) \neq 0$, $m = \pm m^2$. Squaring, $m^2 = m^4$ and dividing by $m$, $m = m^3$. Since $f(m^3) = 0$, $f(m) = 0$, which is a contradiction for $m \neq 1$. However, if we plug in $x = 1$ with $f(1) = 1$ and $y$ as an arbitrary large number with $f(y) = 0$ into the original equation, we get $0 = 1$ which is a clear contradiction, so our only solutions are $f(x) = 0$ and $f(x) = x^2$.

 



\textbf{Solution 2}

Given that the range of f consists entirely of integers, it is clear that the LHS must be an integer and $f(yf(y))$ must also be an integer, therefore $\frac{f(x)^2}{x}$ is an integer. If $x$ divides $f(x)^2$ for all integers $x \ne 0$, then $x$ must be a factor of $f(x)$, therefore $f(0)=0$. Now, by setting $y=0$ in the original equation, this simplifies to $xf(-x)=\frac{f(x)^2}{x}$. Assuming $x \ne 0$, we have $x^2f(-x)=f(x)^2$. Substituting in $-x$ for $x$ gives us $x^2f(x)=f(-x)^2$. Substituting in $\frac{f(x)^2}{x^2}$ in for $f(-x)$ in the second equation gives us $x^2f(x)=\frac{f(x)^4}{x^4}$, so $x^6f(x)=f(x)^4$. In particular, if $f(x) \ne 0$, then we have $f(x)^3=x^6$, therefore $f(x)$ is equivalent to $0$ or $x^2$ for every integer $x$. Now, we shall prove that if for some integer $t \ne 0$, if $f(t)=0$, then $f(x)=0$ for all integers $x$. If we assume $f(y)=0$ and $y \ne 0$ in the original equation, this simplifies to $xf(-x)+y^2f(2x)=\frac{f(x)^2}{x}$. However, since $x^2f(-x)=f(x)^2$, we can rewrite this equation as $\frac{f(x)^2}{x}+y^2f(2x)=\frac{f(x)^2}{x}$, $y^2f(2x)$ must therefore be equivalent to $0$. Since, by our initial assumption, $y \ne 0$, this means that $f(2x)=0$, so, if for some integer $y \ne 0$, $f(y)=0$, then $f(x)=0$ for all integers $x$. The contrapositive must also be true, i.e. If $f(x) \ne 0$ for all integers $x$, then there is no integral value of $y \ne 0$ such that $f(y)=0$, therefore $f(x)$ must be equivalent for $x^2$ for every integer $x$, including $0$, since $f(0)=0$. Thus, $f(x)=0, x^2$ are the only possible solutions.

 



\textbf{Solution 3}

Let's assume $f(0)\neq 0.$ Substitute $(x,y)=(2f(0),0)$ to get 
 \[2f(0)^2=f(2f(0))^2/2f(0)+f(0)\]
 \[2f(0)^2(2f(0)-1)=f(2f(0))^2\]

 This means that $2(2f(0)-1)$ is a perfect square. However, this is impossible, as it is equivalent to $2\pmod{4}.$ Therefore, $f(0)=0.$ Now substitute $x\neq 0, y=0$ to get 
 \[xf(-x)=\frac{f(x)^2}{x} \implies x^2f(-x)=f(x)^2.\]Similarly, 
 \[x^2f(x)=f(-x)^2.\]From these two equations, we can find either $f(x)=f(-x)=0,$ or $f(x)=f(-x)=x^2.$ Both of these are valid solutions on their own, so let's see if there are any solutions combining the two. 

 Let's say we can find $f(x)=x^2, f(y)=0,$ and $x,y\neq 0.$ Then 
 \[xf(-x)+y^2f(2x)=f(x)^2/x.\]
 \[y^2f(2x)=x-x^3.\] (NEEDS FIXING: $f(x)^2/x= x^4/x = x^3$, so the RHS is $0$ instead of $x-x^3$.)

 If $f(2x)=4x^2,$ then $y^2=\frac{x-x^3}{4x^2}=\frac{1-x^2}{4x},$ which is only possible when $y=0.$ This contradicts our assumption. Therefore, $f(2x)=0.$ This forces $x=\pm 1$ due to the right side of the equation. Let's consider the possibility $f(2)=0, f(1)=1.$ Substituting $(x,y)=(2,1)$ into the original equation yields  
 \[0=2f(0)+1f(2)=0+f(1)=1,\] which is impossible. So $f(2)=f(-2)=4$ and there are no solutions "combining" $f(x)=x^2$ and $f(x)=0.$

 Therefore our only solutions are $\boxed{f(x)=0}$ and $\boxed{f(x)=x^2.}$

 



\textbf{Solution 4}

Let the given assertion be $P(x, y)$. We try $P(x, 0)$ and get $xf(2c-x)=f(x)^2/x+c$, where $f(0)=c$. We plug in $x=c$ and get $cf(c)=f(c)^2/c+c$. Rearranging and solving for $c^2$ gives us $c^2=\frac{f(c)^2}{f(c)-1}$. Obviously, the only $c$ that works such that the RHS is an integer is $c=0$, and thus $f(0)=0$.

 We use this information on assertion $P(x,0)$ and obtain $xf(-x)=f(x^2)/x$, or $f(-x)=\frac{f(x)^2}{x^2}$. Thus, $f(x)$ is an even function. It follows that $f(x)=0, x^2$ for each $x$. We now prove that $f(x)=x^2$, f(x)=0$ are the only solutions. [in progress] ~SigmaPiE

 


	\item (\textit{USAJMO 2014}) Let $b\geq 2$ be an integer, and let $s_b(n)$ denote the sum of the digits of $n$ when it is written in base $b$.  Show that there are infinitely many positive integers that cannot be represented in the form $n+s_b(n)$, where $n$ is a positive integer.


 



\textbf{Solution}

Define $S(n) = n + s_b(n)$, and call a number \textit{unrepresentable} if it cannot equal $S(n)$ for a positive integer $n$.We claim that in the interval $(b^p, b^{p+1}]$ there exists an unrepresentable number, for every positive integer $p$.

 If $b^{p+1}$ is unrepresentable, we're done. Otherwise, time for our lemma:

 Lemma: Define the function $f(p)$ to equal the number of integers x less than $b^p$ such that $S(x) \ge b^p$. If $b^{p+1} = S(y)$ for some y, then $f(p+1) > f(p)$.

 Proof: Let $F(p)$ be the set of integers x less than $b^p$ such that $S(x) \ge b^p$. Then for every integer in $F(p)$, append the digit $(b-1)$ to the front of it to create a valid integer in $F(p+1)$. Also, notice that $(b-1) \cdot b^p \le y < b^{p+1}$. Removing the digit $(b-1)$ from the front of y creates a number that is not in $F(p)$. Hence, $F(p) \rightarrow F(p+1)$, but there exists an element of $F(p+1)$ not corresponding with $F(p)$, so $f(p+1) > f(p)$.

 Note that our lemma combined with the Pigeonhole Principle essentially proves the claim. Therefore, because there are infinitely many disjoint intervals containing an unrepresentable number, there are infinitely many unrepresentable numbers.

 



\textbf{Solution 2 (Simple)}

Let $f(n) = n + s_b(n)$. It is easy to see that $f(b^2) = b^2 + 1$, and $f(b^2 - b + 1) = b^2 - b + 1 + b - 1 + 1 = b^2 + 1$, because the digits of $b^2 - b + 1$ are $b-1$ and $1$. Additionally, if we take any number, let's say $x$, and multiply it by $b^3$, and then add it to $b^2$ and $b^2 - b + 1$, we still obtain $f(b^3x + b^2) = f(b^3x + b^2 - b + 1)$, because we are simply adding $b^3x + s_b(x)$ to each value. 

 Thus, we conclude that there are infinitely many numbers that are counted twice in $f(n)$, and that these numbers come $b^3$ apart. Finally, it is clear that for a large $N$, the number of numbers that are counted at least twice is at least $\left\lfloor\frac{N}{b^3}\right\rfloor$. Because $f(n) > n$ for all $n$, this means that there are at least $\left\lfloor\frac{N}{b^3}\right\rfloor$ numbers uncounted, and because this is unbounded, the proof is complete.

 \textbf{Operation Intuitive Number Theory (INT)- Solution 2} I hope this solution is quite intuitive, because it is without complicated notation. It didn't take me very long to discover.As with Solution 1, define $S(n)=n+s_b(n)$.We will use numbers from $b^k$ to $b^{k+1}-1$ for induction. Call this interval Class $k$.

 Start with $b^0=1$ and go up to $b-1$. There are $b-1$ numbers covered. Easily $S(n)$ ranges from $1+1=2$ to $b-1+b-1=2b-2$, or a range of $2b-3$. Thus, there are at LEAST $b-2$ numbers lacking coverage.Now, in order to fill up these gaps, we consult Class $1$: $b^1=b$ to $b^2-1$. If we are to fill up all these gaps, then we need at least $b-2$ numbers in the next series with their $S$ values. Unfortunately, there are only at MOST $b-3$ numbers that can satisfy a value, from $b$ to $2b-3$, otherwise the $S$ value is too big (note that none of them are zero)! Thus, between Class 0 and 1, there is at least one value lacking.

 After setting up the base case, consider the Class $k$ numbers from $b^k$ to $b^{k+1}-1$ there are clearly $b^{k+1}-b^k$ integers. Yet the $S(n)$ can range from $b^k+1=b^k+1$ to $(b^{k+1}-1)+(k+1)(b-1)=b^{k+1}+b(k+1)-(k+1)$. This gives a range of $b^{k+1}-b^k+b(k+1)-(k+1)$. This leaves an extra $b(k+1)-(k+1)$ numbers. Now, we invoke the next class: $k+1$. Again, to fill in the gap there are sadly only $b(k+1)-(k+2)$ numbers available: $b^{k+1}$ to $b^{k+1}+b(k+1)-(k+3)$, because $s_b(n)$ is at least 1.

 By induction we are done. Because each Class $k$ for integral $k$ at least $0$ misses at least one $S$ value, we miss an infinite number of numbers. Game over! 

 Note: If you are unconvinced of the range of $S(n)$ values we go from the smallest value of both the number and the $s_b(n)$, or $b^k$ and $1$, to the greatest of both, i.e. $b^{k+1}-1$ and $k(b+1)$, because in that series the largest possible sum is a bunch of $b-1$s.

 


	\item (\textit{USAJMO 2014}) Let $k$ be a positive integer. Two players $A$ and $B$ play a game on an infinite grid of regular hexagons. Initially all the grid cells are empty. Then the players alternately take turns with $A$ moving first. In his move, $A$ may choose two adjacent hexagons in the grid which are empty and place a counter in both of them. In his move, $B$ may choose any counter on the board and remove it. If at any time there are $k$ consecutive grid cells in a line all of which contain a counter, $A$ wins. Find the minimum value of $k$ for which $A$ cannot win in a finite number of moves, or prove that no such minimum value exists.


 



\textbf{Solution}

The answer is $k=6$. We prove that $A$ can win for $k=5$ (which hence proves it for $k<5$ as well) and show that $B$ can thwart $A$ for $k\geq 6$.

 Arrange the board so that a pair of opposite sides are horizontal. Create a coordinate system on the board by setting the center of some hexagon as the origin and setting the hexagons directly above and above-and-right as $(0, 1)$ and $(1, 0)$, respectively. Then, for example, the below-and-right hexagon touching the origin is $(1, -1)$. So two hexagons touch if their coordinate difference is one of these.

 Now for $k=5$, person $A$ places his counters only in $\{0, 1, 2, 3, 4\}\times\{0, 1\}$. Note that if, at $A$'s turn, there are 4 counters in either column, then $A$ can win immediately, so let us assume that in both columns there are at least 2 missing, meaning that at most $6$ counters are on the board. We would like to find when $A$ cannot play under these circumstances. If we look at the disjoint sets

 $\{(0, 0), (0, 1), (1, 0)\}, \{(1, 1), (2, 0), (2, 1)\}, \{(3, 0), (3, 1), (4, 0)\}, \{(4, 1)\}$

 we see that either some set has at least $2$ hexagons without counters, in which case $A$ can move, or all four sets have exactly $1$ missing counter. Similarly for the sets

 $\{(0, 0)\}, \{(0, 1), (1, 0), (1, 1)\}, \{(2, 0), (2, 1), (3, 0)\}, \{(3, 1), (4, 0), (4, 1)\}$

 So both $(0, 0), (4, 1)$ have no token on them. This means that $(0, 1), (1, 0), (3, 1), (4, 0)$ do. Thus $(3, 0)$ and $(1, 1)$ do not. So this is the only situation in which $A$ can neither win immediately nor play in only these 10 hexagons.

 So $A$ plays only in these 10 hexagons until either he has a win or he can't anymore. If he wins, then we're done. Otherwise, $A$ plays in hexagons $(5, 0)$ and $(5, 1)$. Then $B$ either removes $(5, 1)$ so that $A$ can win at $(2, 0)$, or $B$ removes $(5, 0)$ and $A$ plays at $(5, 0)$ and $(4, 1)$ and then at either $(3, 0)$ or $(1, 1)$ the next turn, or $B$ removes one of $\{0, 1, 2, 3, 4\}\times\{0, 1\}$ in which case $A$ can either win immediately at $(3, 0)$ or can play in both columns, and then win the next turn.

 Now if $k\geq 6$, then if $A$ plays on anything in the lattice generated by $(2, -1)$ and $(1, 1)$, that is $(2a+b, b-a)$ for $a, b$ integers, then $B$ removes it. Otherwise, $B$ removes any of $A$'s counters. This works because in order for $A$ to win, there must be at least 2 counters in this lattice, but $A$ can only put a counter on $1$ at any time, so there's at most 1 on the lattice at any time.

 So $A$ wins if $k\leq 5$, and $B$ wins if $k\geq 6$.

 


	\item (\textit{USAJMO 2014}) Let $ABC$ be a triangle with incenter $I$, incircle $\gamma$ and circumcircle $\Gamma$.  Let $M,N,P$ be the midpoints of sides $\overline{BC}$, $\overline{CA}$, $\overline{AB}$ and let $E,F$ be the tangency points of $\gamma$ with $\overline{CA}$ and $\overline{AB}$, respectively.  Let $U,V$ be the intersections of line $EF$ with line $MN$ and line $MP$, respectively, and let $X$ be the midpoint of arc $BAC$ of $\Gamma$.


 (a) Prove that $I$ lies on ray $CV$.


 (b) Prove that line $XI$ bisects $\overline{UV}$.


 



\textbf{Solution 1}

(a) Let $V' = \overline{CI} \cap \overline{EF}$. We claim $BIV'F$ is cyclic. This follows by Angle chasing: 
 \[\measuredangle V'FB \equiv \measuredangle V'FE \equiv 90^\circ -\frac{A}{2} \pmod{180 ^\circ} \text{ and } \measuredangle V'IB \equiv \frac{B}{2}+\frac{C}{2} \equiv 90^\circ -\frac{A}{2} \pmod{180^\circ}\] Thus, $\measuredangle BV'C = \measuredangle BV'I = \measuredangle BFI = 90^\circ$. It follows that $(BV'C)$ is a circle centered at $M$, so $\angle V'MB = 2\angle V'CB = C$.

 But, $\angle PMC = C$, so $V' \in \overline{MP}$ and thus $V' = V$, as needed.

 
(b) Apply a homothety at $I$ sending $\overline{XI} \cap \overline{EF}$ to $X$. Notice that $\overline{EF} \parallel \overline{AX}$, so the homothety sends $\overleftrightarrow{EF}$ to $\overleftrightarrow{AX}$. 

 Thus, $U' = \overline{AX}\cap \overline{BI} = I_B$ and $V'= \overline{AX}\cap \overline{CI} = I_C$, where $I_B, I_C$ are the respective excenters. 

 By duality, $A, B, C$ are the feet of the altitude of $\triangle I_AI_BI_C$, so $\Gamma$ is the nine-point circle of $I_AI_BI_C$. 

 But, $X = \overline{I_BI_C} \cap \Gamma$ other than $A$, so $X$ must be the midpoint of $I_BI_C$, as needed.

 
Note 1: In (a), we use directed angles with half-angles, which typically doesn't work. I circumvent this issue by considering the half-angles themselves mod 180, and this means there is no need to half any of the angles.

 Note 2: In (b), you can prove more generally that $B, C, I_B, I_C$ sit on a circle centered at $X$. Maybe call it Fact 6.

 
-sillybone

 


 



\textbf{Solution 2}


                    \begin{figure}[H]
                        \centering
                        \includegraphics[width=0.5\textwidth]{Solution/USAJMO_I_2014_6_Solution_Figure_1.png}
                    \end{figure}
                    

 (a)

 \textbf{Solution 1:} We will prove this via contradiction: assume that line $IC$ intersects line $MP$ at $Q$ and line $EF$ and $R$, with $R$ and $Q$ not equal to $V$. Let $x = \angle A/2 = \angle IAE$ and $y = \angle C/2 = \angle ICA$. We know that $\overline{MP} \parallel \overline{AC}$ because $MP$ is a midsegment of triangle $ABC$; thus, by alternate interior angles (A.I.A) $\angle MVE = \angle FEA = (180^\circ - 2x) / 2 = 90^\circ - x$, because triangle $AFE$ is isosceles. Also by A.I.A, $\angle MQC = \angle QCA = y$. Furthermore, because $AI$ is an angle bisector of triangle $AFE$, it is also an altitude of the triangle; combining this with $\angle QIA = x + y$ from the Exterior Angle Theorem gives $\angle FRC = 90 - x - y$. Also, $\angle VRQ = \angle FRC = 90 - x - y$ because they are vertical angles. This completes part (a).

 \textbf{Solution 2:} First we show that the intersection $V'$of $MP$ with the internal angle bisector of $C$ is the same as the intersection $V''$ of $EF$ with the internal angle bisector of $C.$ Let $D$ denote the intersection of $AB$ with the internal angle bisector of $C,$ and let $a,b,c$ denote the side lengths of $BC, AC, AB.$ By Menelaus on $V'', F, E,$ with respect to $\triangle ADC,$ 
 \[\frac{AF}{FD}\cdot \frac{DV''}{V''C}\cdot \frac{CE}{EA}=-1\]
 \[\frac{DV''}{V''C}=-\frac{\frac{bc}{a+b}+\frac{a-b-c}{2}}{\frac{a+b-c}{2}}=\frac{b-a}{b+a}.\] Similarly, 
 \[\frac{BP}{PD}\cdot \frac{DV'}{V'C}\cdot \frac{CM}{MB}=-1\] 
 \[\frac{DV'}{V'C}=-\frac{\frac{ac}{a+b}-\frac{c}{2}}{\frac{c}{2}}=\frac{b-a}{b+a}.\] Since $V'$ and $V''$ divide $CD$ in the same ratio, they must be the same point. Now, since $\frac{b-a}{b+a}>-1,$ $I$ lies on ray $CV.$ $\blacksquare$

 \textbf{Solution 3:} By the Iran Lemma, we know $CI, EF, MP$ concur, so Part $\text{A}$ follows easily.

 (b)

 \textbf{Solution 1:} Using a similar argument to part (a), point U lies on line $BI$. Because $\angle MVC = \angle VCA = \angle MCV$, triangle $VMC$ is isosceles. Similarly, triangle $BMU$ is isosceles, from which we derive that $VM = MC = MB = MU$. Hence, triangle $VUM$ is isosceles.

 Note that $X$ lies on both the circumcircle and the perpendicular bisector of segment $BC$. Let $D$ be the midpoint of $UV$; our goal is to prove that points $X$, $D$, and $I$ are collinear, which equates to proving $X$ lies on ray $ID$.

 Because $MD$ is also an altitude of triangle $MVU$, and $MD$ and $IA$ are both perpendicular to $EF$, $\overline{MD} \parallel \overline{IA}$. Furthermore, we have $\angle VMD = \angle UMD = x$ because $APMN$ is a parallelogram. (incomplete)

 \textbf{Solution 2:} Let $I_A$, $I_B$, and $I_C$ be the excenters of $ABC$. Note that the circumcircle of $ABC$ is the nine-point circle of $I_AI_BI_C$. Since $AX$ is the external angle bisector of $\angle BAC$, $X$ is the midpoint of $I_BI_C$. Now $UV$ and $I_BI_C$ are parallel since both are perpendicular to the internal angle bisector of $\angle BAC$. Since $IX$ bisects $I_BI_C$, it bisects $UV$ as well.

 \textbf{Solution 3:} Let $X'$ be the antipode of $X$ with respect to the circumcircle of triangle $ABC$. Then, by the Incenter-Excenter lemma, $X'$ is the center of a circle containing $B$, $I$, and $C$. Because $XX'$ is a diameter, $XB$ and $XC$ are tangent to the aforementioned circle; thus by a well-known symmedian lemma (the Tangent Construction Lemma), $XI$ is the $I$-symmedian of triangle $IBC$. From part (a); we know that $BVUC$ is cyclic with diameter $BC$ (we can derive a similar argument for point $U$); thus, from similar triangles, we know that $XI$ is also the median of triangle $VIU$, and we are done.

 \textbf{Solution 4:} Let $M_b, M_c$ be the midpoints of arcs $CA, AB$ respectively, and let $T_a$ be the tangency point between the $A$-mixtilinear incircle of $ABC$ and $\Gamma$. It's well-known that $T_a \in XI$, $T_aX$ bisects $M_bM_c$, $AI \perp EF$, and $AI \perp M_bM_c$.

 Now, it's easy to see $EF \parallel M_bM_c$, so $IUV$ and $IM_bM_c$ are homothetic at $I$. But $T_aX \equiv IX$ bisects $M_bM_c$, so Part $\text{B}$ follows directly from the homothety. ~ ike.chen

 \textbf{Another Solution for Part (a)} Let $\angle CAI = \angle BAI = x,$ $\angle ACI = \angle BCI = z,$ and $\angle ABI = \angle CBI = y$. In this solution, we use degrees.

 
Since $AE = AF$ (they are both tangents), we know that $\Delta AEF$ is isosceles with vertex angle bisected by $\overline{AI}$. Therefore, $\overline{AI} \perp \overline{EF}$.This means that $\angle AEF = 90 - x$, so $\angle CEF = 90 + x$.

 
Let $V$ be the intersection of $\overleftrightarrow{EF}$ and $\overleftrightarrow{CI}$, so
 \[\angle EVC = 180 - (90 + x) - z = 90 - x - z = \frac{180-2x-2z}{2}= \frac{2y}{2} = y,\] so
 \[\angle FVI = 180 - y = 180 - \angle FBI.\]Therefore, $BFVI$ is cyclic, meaning
 \[\angle BVC = \angle BVI = \angle BFI = 90.\]This means that $M$ is the circumcenter of $\Delta BVC$, so $MV = CM$, so $\angle CVM = z \rightarrow \angle CMV = 180 - 2z$.

 
Then $\angle CMV$ is supplementary to $\angle ACB$, so $\overleftrightarrow{VM} \parallel \overleftrightarrow{AC} \parallel \overleftrightarrow{MP}$, so $V$ is on $\overleftrightarrow{MP}$, so the three lines $\overleftrightarrow{EF}, \overleftrightarrow{CI}, \overleftrightarrow{MP}$ are concurrent at $V$, meaning that $I$ is on $\overrightarrow{CI}$, as desired.

 


\end{enumerate}